\section{Detailed derivation of the conjugate force and the response function of some observables}
\subsection{General formula of the conjugate force}
\quad From the effective potential in (\ref{steady-state prob}), the conjugate force is defined as $\vec{\mathscr{F}} \equiv -\nabla_p \Phi$. Then $\nabla_p \ln{\psi_{ss}} = \delta \vec{\mathscr{F}}$. Since the effective potential $\Phi(\vec{x}; \vec{p}, \vec{X}(\vec{p}))$ depends on the fixed point and the fixed point itself depends on $\vec{p}$, one obtains
\begin{equation}
    \mathscr{F}_\alpha = -\frac{\partial \Phi}{\partial p_\alpha} - \sum_s \frac{\partial \Phi}{\partial X_s} \frac{\partial X_s}{\partial p_\alpha} = -\frac{\partial \Phi}{\partial p_\alpha} + \sum_s \frac{\partial \Phi}{\partial X_s} \left[ \mathbf{Q}^{-1} \nabla_p \vec{F} \right]_{s\alpha}.
    \label{scriptF}
\end{equation} 
Then the dependence of $\Phi$ on the fixed point is
\begin{eqnarray}
    \frac{\partial \Phi}{\partial X_s} &=& \frac{\partial }{\partial X_s} \left( \frac{1}{2} \sum_{\ell,m} ({\bf K}_0^{-1})_{\ell m} \delta x_\ell \delta x_m \right)\\
    &=&-\sum_{m}({\bf K}_0^{-1})_{sm}\delta x_m + \frac{1}{2}\sum_{\ell,m} \sum_{\mu,\nu} \frac{\partial ({\bf K}_0^{-1})_{\ell m}}{\partial Q_{\mu\nu}}\frac{\partial Q_{\mu\nu}}{\partial X_s} \delta x_\ell \delta x_m\nonumber\\
    &=& -\sum_{m}({\bf K}_0^{-1})_{sm}\delta x_m +\frac{1}{2}\sum_{\ell,m} \Bigg\{\frac{\partial ({\bf K}_0^{-1})_{\ell m}}{\partial Q_{ss}} \left[f''_s(X_s;r_s)+\sum_{n\neq s} W_{s n}\partial^2_{1}h(X_s,X_n)\right]\nonumber\\ &+&\sum_{\mu}\frac{\partial ({\bf K}_0^{-1})_{\ell m}}{\partial Q_{\mu\mu}} W_{\mu s} \partial_{1}\partial_2 h(X_\mu,X_s) + \sum_{\nu}\frac{\partial ({\bf K}_0^{-1})_{\ell m}}{\partial Q_{s\nu}} W_{s \nu}  \partial_{1}\partial_2 h(X_s,X_\nu)\nonumber\\
    &+&\sum_{\mu} \frac{\partial ({\bf K}_0^{-1})_{\ell m}}{\partial Q_{\mu s} }W_{\mu s} \partial^2_{2}h(X_\mu,X_s) \Bigg\} \delta x_\ell \delta x_m\nonumber,
    \label{pphipX}
\end{eqnarray}
where we use
\begin{eqnarray}
    \frac{\partial Q_{\mu\nu}}{\partial X_s} &=& \delta_{\mu s} \delta_{\mu\nu} \left[ f''_s(X_s;r_s)+\sum_{m\neq s} W_{sm} \partial_1^2 h(X_s,X_m) \right] + \delta_{\mu\nu} W_{\mu s} \partial_1\partial_2 h(X_\mu,X_s)\nonumber\\
    &+&\delta_{\mu s} W_{s\nu} \partial_1 \partial_2 h (X_s,X_\nu)+\delta_{\nu s} W_{\mu s} \partial_2^2 h (X_\mu,X_s).
\end{eqnarray}
$ \frac{\partial {\bf K}_0^{-1}}{\partial Q_{\mu\nu}}  $ denotes the matrix whose $i,j$ elements are $\frac{\partial ({\bf K}_0^{-1})_{ij}}{\partial Q_{\mu\nu}}$.\par
In the linearized noisy network dynamics, the steady-state probability distribution is approximated by a Gaussian distribution near its mean $\langle \vec{x} \rangle$, with $\Phi = \frac{1}{2} \delta \vec{x}^\intercal \mathbf{K}^{-1}_0 \delta \vec{x} = \frac{1}{2} (\vec{x}-\langle \vec{x} \rangle)^\intercal \mathbf{K}^{-1}_0 (\vec{x}-\langle \vec{x} \rangle)$. Note that $\langle {\vec x} \rangle$ depends on the parameters ${\vec p}$ through $\psi_{ss}$, and this dependence is generally difficult to compute explicitly in high-dimensional nonlinear Langevin systems. We therefore approximate $\langle \vec{x} \rangle$ by $\langle \vec{X} \rangle$. Nevertheless, the conjugate force ${\vec{\mathscr F}}$ and correlations such as $\langle\delta \mathbf{B} \delta\vec{\mathscr F}\rangle_{ss,{\vec p}}$ can still be computed explicitly from (\ref{scriptF}) using the functional dependence of ${\bf K}_0^{-1}$ on ${\vec p}$. From the Lyapunov equation (\ref{Lyapunov_eq}), one sees that ${\bf K}_0^{-1}={\bf K}_0^{-1}({\bf Q}({\vec p},{\vec X}({\vec p})))$. Hence the conjugate force corresponding to the parameter $p_\alpha$ can be written as
\begin{eqnarray}
{\mathscr F}_\alpha&=& - \frac{1}{2} \sum_{\ell, m}\sum_{\mu,\nu} \frac{\partial ({\bf K}_0^{-1})_{\ell m}} {\partial Q_{\mu\nu}}\frac {\partial Q_{\mu\nu}}{\partial p_\alpha} \delta x_\ell \delta x_m -\sum_{\ell, m} \left({\bf Q}^{-1} \nabla_p{\vec { F}} \Big|_{\vec X}\right)_{\ell\alpha}({\bf K}_0^{-1})_{\ell m}\delta x_m \label{Falp}\\
&+&\frac{1}{2} \sum_s \left({\bf Q}^{-1} \nabla_p{\vec { F}} \Big|_{\vec X}\right)_{s\alpha} \sum_{\ell,m} \Bigg\{ \frac{\partial (\mathbf{K}_0^{-1})_{\ell m}}{\partial Q_{ss}}\left( f_s''(X_s;r_s)+\sum_{n\neq s} W_{sn} \partial^2_1 h(X_s, X_n) \right)\nonumber\\
&+& \sum_\mu \frac{\partial (\mathbf{K}_0^{-1})_{\ell m}}{\partial Q_{\mu\mu}} W_{\mu s} \partial_1 \partial_2 h(X_\mu, X_s) + \sum_\nu \frac{\partial (\mathbf{K}_0^{-1})_{\ell m}}{\partial Q_{s\nu}} W_{s\nu}\partial_1\partial_2 h(X_s,X_\nu)\nonumber\\ &+& \sum_\mu \frac{\partial (\mathbf{K}_0^{-1})_{\ell m}}{\partial Q_{\mu s}} W_{\mu s} \partial^2_2 h(X_\mu, X_s) \Bigg\} \delta x_\ell \delta x_m .\nonumber
\end{eqnarray}
One can derive, or directly verify from (\ref{Falp}), that
\begin{equation}
    \frac{\partial X_i}{\partial p_\alpha} =\langle \delta x_i {\mathscr F}_\alpha\rangle,
\end{equation}
where $\langle \delta x_i \delta x_\ell \delta x_m \rangle=0$. By using (\ref{inverse_partialKo}), we can write down the ensemble average of the conjugate force
\begin{eqnarray}
\langle{{\mathscr F}_\alpha}\rangle&=& \frac{1}{2} \sum_{\mu,\nu} \frac {\partial Q_{\mu\nu}}{\partial p_\alpha} \text{Tr}\left({\bf K}_0^{-1} \frac{\partial {\bf K}_0} {\partial Q_{\mu\nu}}\right)-\frac{1}{2} \sum_s \left({\bf Q}^{-1} \nabla_p{\vec { F}} \Big|_{\vec X}\right)_{s\alpha}\Bigg\{ \Bigg[ f_s''(X_s;r_s)\\
&+&\sum_{n\neq s} W_{sn} \partial^2_1 h(X_s, X_n) \Bigg] \text{Tr}\left({\bf K}_0^{-1} \frac{\partial {\bf K}_0} {\partial Q_{ss}}\right)+\sum_\mu W_{\mu s} \partial_1 \partial_2 h(X_\mu, X_s) \text{Tr}\left({\bf K}_0^{-1} \frac{\partial {\bf K}_0} {\partial Q_{\mu\mu}}\right)\nonumber\\
&+& \sum_\nu W_{s\nu}\partial_1\partial_2 h(X_s,X_\nu) \text{Tr}\left({\bf K}_0^{-1} \frac{\partial {\bf K}_0} {\partial Q_{s\nu}}\right)+ \sum_\mu W_{\mu s} \partial^2_2 h(X_\mu, X_s) \text{Tr}\left({\bf K}_0^{-1} \frac{\partial {\bf K}_0} {\partial Q_{\mu s}}\right) \Bigg\}.\nonumber
\end{eqnarray}
\subsection{Conjugate force with respect to the change in the intrinsic node parameter r}
\quad Consider perturbations of an intrinsic node parameter $r_\alpha$.
\begin{equation}
    \left({\bf Q}^{-1} \nabla_r{\vec {F}}|_{\vec X} \right)_{s\alpha} =
    Q_{s\alpha}^{-1}\frac{\partial f_\alpha}{\partial r_\alpha} \bigg|_{\vec X},\quad
    \frac{\partial Q_{\mu\nu}}{\partial r_\alpha}=\frac{\partial f'_\alpha}{\partial r_\alpha}\bigg|_{\vec X}\delta_{\mu\alpha}\delta_{\mu\nu},\label{QgradF}
\end{equation}
and ${\mathscr F}_\alpha$  can then be calculated from (\ref{Falp}) to give
\begin{eqnarray}
    {\mathscr F}_\alpha&=& -\frac{1}{2} \frac{\partial f'_\alpha}{\partial r_\alpha}\Bigg|_{\vec{X}} \sum_{\ell,m} \frac{\partial (\mathbf{K}_0^{-1})_{\ell m}}{\partial Q_{\alpha\alpha}} \delta x_\ell \delta x_m - \frac{\partial f_\alpha}{\partial r_\alpha}\Bigg|_{\vec{X}} \sum_{\ell,m} Q^{-1}_{\ell\alpha} (\mathbf{K}_0^{-1})_{\ell m} \delta x_m\\
    &+&\frac{1}{2}\frac{\partial f_\alpha}{\partial r_\alpha}\Bigg|_{\vec{X}} \sum_s {Q}^{-1}_{s\alpha} \sum_{\ell,m} \Bigg\{ \frac{\partial (\mathbf{K}_0^{-1})_{\ell m}}{\partial Q_{ss}}\left( f_s''(X_s;r_s)+\sum_{n\neq s} W_{sn} \partial^2_1 h(X_s, X_n) \right)\nonumber\\
    &+& \sum_\mu \frac{\partial (\mathbf{K}_0^{-1})_{\ell m}}{\partial Q_{\mu\mu}} W_{\mu s} \partial_1 \partial_2 h(X_\mu, X_s) + \sum_\nu \frac{\partial (\mathbf{K}_0^{-1})_{\ell m}}{\partial Q_{s\nu}} W_{s\nu}\partial_1\partial_2 h(X_s,X_\nu)\nonumber\\ &+& \sum_\mu \frac{\partial (\mathbf{K}_0^{-1})_{\ell m}}{\partial Q_{\mu s}} W_{\mu s} \partial^2_2 h(X_\mu, X_s) \Bigg\} \delta x_\ell \delta x_m .\nonumber
\end{eqnarray}
Average quantities involving ${\mathscr F}_\alpha$ can then be calculated. For example, one obtains
\begin{eqnarray}
\langle{{\mathscr F}_\alpha}\rangle&=& \frac{1}{2} \frac{\partial f'_\alpha}{\partial r_\alpha}\bigg|_{\vec X} \text{Tr}\left({\bf K}_0^{-1} \frac{\partial {\bf K}_0} {\partial Q_{\alpha\alpha}}\right)-\frac{1}{2}\frac{\partial f_\alpha}{\partial r_\alpha} \bigg|_{\vec X} \sum_s Q_{s\alpha}^{-1}\Bigg\{ \Bigg[ f_s''(X_s;r_s)\\
&+&\sum_{n\neq s} W_{sn} \partial^2_1 h(X_s, X_n) \Bigg] \text{Tr}\left({\bf K}_0^{-1} \frac{\partial {\bf K}_0} {\partial Q_{ss}}\right)+\sum_\mu W_{\mu s} \partial_1 \partial_2 h(X_\mu, X_s) \text{Tr}\left({\bf K}_0^{-1} \frac{\partial {\bf K}_0} {\partial Q_{\mu\mu}}\right)\nonumber\\
&+& \sum_\nu W_{s\nu}\partial_1\partial_2 h(X_s,X_\nu) \text{Tr}\left({\bf K}_0^{-1} \frac{\partial {\bf K}_0} {\partial Q_{s\nu}}\right)+ \sum_\mu W_{\mu s} \partial^2_2 h(X_\mu, X_s) \text{Tr}\left({\bf K}_0^{-1} \frac{\partial {\bf K}_0} {\partial Q_{\mu s}}\right) \Bigg\},\nonumber
\end{eqnarray}
To calculate the response of observables that depend on $\delta \vec{x}$, one must consider multipoint correlations between $\delta x$ and $\mathscr{F}_{\alpha}$. Higher-order correlations of $\delta x_i$ can be evaluated using the cumulant expansion for a high-dimensional Gaussian distribution. In this way, quantities such as $\langle \delta x_i \delta x_j \cdots \mathscr{F}_{\alpha}\rangle$ can be computed from the corresponding moments of $\delta x$. Here we display only the two basic results, $\langle \delta x_i\mathscr{F}_\alpha\rangle$ and $\langle \delta x_i \delta x_j{\mathscr F}_\alpha\rangle$, namely
\begin{equation}
    \langle \delta x_i \mathscr{F}_{\alpha}\rangle = \frac{\partial X_i}{\partial r_\alpha},
    \label{dxdr}
\end{equation}
\begin{eqnarray}
\langle \delta x_i \delta x_j{\mathscr F}_\alpha\rangle&=& \frac{1}{2}\frac{\partial f'_\alpha}{\partial r_\alpha}\bigg|_{\vec X}\left[ ({\bf K}_0)_{ij}\text{Tr}\left({\bf K}^{-1}_0\frac{\partial {\bf K}_0}{\partial Q_{\alpha \alpha}} \right)+2 \left(\frac{\partial {\bf K}_0}{\partial Q_{\alpha \alpha}}\right)_{ij} \right]\label{dxdxscrF}\\
&-&\frac{1}{2} \frac{\partial f_\alpha}{\partial r_\alpha}\bigg|_{\vec X} \sum_s Q_{s\alpha}^{-1}\Biggl\{ \left[ f_s''(X_s;r_s)+  \sum_n W_{sn}\partial_1^2h(X_s,X_n)\right] \times\nonumber\\
& &\left[ ({\bf K}_0)_{ij} \text{Tr}\left({\bf K}^{-1}_0\frac{\partial {\bf K}_0}{\partial Q_{ss}}\right)+2 \left(\frac{\partial {\bf K}_0}{\partial Q_{ss}} \right)_{ij}\right] \nonumber\\
&+&\sum_\mu W_{\mu s}\partial_1\partial_2h(X_\mu,X_s) \left[ ({\bf K}_0)_{ij} \text{Tr} \left({\bf K}_0^{-1}\frac{\partial {\bf K}_0}{\partial Q_{\mu\mu}} \right)+2 \left(\frac{\partial {\bf K}_0}{\partial Q_{\mu\mu}} \right)_{ij}\right] \nonumber\\
&+& \sum_\nu W_{s\nu} \partial_1\partial_2h(X_s,X_\nu)  \left[ ({\bf K}_0)_{ij}\text{Tr} \left({\bf K}_0^{-1}\frac{\partial {\bf K}_0}{\partial Q_{s\nu}} \right)+2 \left(\frac{\partial {\bf K}_0}{\partial Q_{s\nu}} \right)_{ij}\right]    \nonumber\\
& +&\sum_\mu W_{\mu s}  \partial_2^2 h(X_\mu,X_s)\left[ ({\bf K}_0)_{ij} \text{Tr} \left({\bf K}_0^{-1} \frac{\partial {\bf K}_0}{\partial Q_{\mu s}} \right)+2 \left(\frac{\partial {\bf K}_0}{\partial Q_{\mu s}} \right)_{ij}\right]\Biggr\}.\nonumber
\end{eqnarray}
\subsection{Conjugate force with respect to the change in the connection weight W}
\quad We next consider perturbations of a network connection $W_{\alpha\beta}$, which directly modify the interaction topology. As in the previous subsection, we first compute the derivatives with respect to the directed connection $W_{\alpha\beta}$:
\begin{eqnarray}
\left({\bf Q}^{-1} \nabla_W{\vec { F}} |_{\vec X}\right)_{s\alpha\beta} &=& 
Q_{s\alpha}^{-1}h(X_\alpha, X_\beta)\label{QgradWF}\\
\frac{\partial Q_{\mu\nu}}{\partial W{_{\alpha\beta}}}&=& \partial_1 h(X_\alpha, X_\beta)\delta_{\mu\nu}\delta_{\mu\alpha}+\partial_2h(X_\alpha, X_\beta)\delta_{\mu\alpha}\delta_{\nu\beta}.\label{dQdW}
\end{eqnarray}
$\boldsymbol{{\mathscr F}} \equiv -\nabla_{\mathbf{W}} \Phi$ is the conjugate force associated with a change in the connection weight $W_{\alpha\beta}$, and it can be calculated from (\ref{Falp}) to give
\begin{eqnarray}
{\mathscr F}_{\alpha\beta}&=& -\frac{1}{2}\sum_{\ell,m} \biggl\{\frac{\partial ({\bf K}_0^{-1})_{\ell m}}{\partial Q_{\alpha \alpha}}\partial_1 h(X_\alpha, X_\beta)+ \frac{\partial ({\bf K}_0^{-1})_{\ell m}}{\partial Q_{\alpha \beta}}\partial_2h(X_\alpha, X_\beta)\biggr\} \delta x_\ell \delta x_m \nonumber\\
&-& h(X_\alpha, X_\beta)\sum_{\ell,m} Q^{-1}_{\ell\alpha}({\bf K}_0^{-1})_{\ell m}\delta x_m\\
&+&\frac{ h(X_\alpha, X_\beta)}{2} \sum_s Q_{s\alpha}^{-1}\sum_{\ell,m}  \Biggl\{ \frac{\partial ({\bf K}_0^{-1})_{\ell m}} {\partial Q_{ss}}\biggl[f_s''+\sum_n W_{sn}  \partial^2_1h(X_s,X_n)\biggr]\nonumber\\
&+& \sum_\mu W_{\mu s} \partial_1\partial_2h(X_\mu,X_s)\frac{\partial ({\bf K}_0^{-1})_{\ell m}} {\partial Q_{\mu\mu}} + \sum_\nu W_{s\nu}\partial_1\partial_2h(X_s,X_\nu)\frac{\partial ({\bf K}_0^{-1})_{\ell m}}{\partial Q_{s\nu}}\nonumber\\
&+& \sum_\mu  W_{\mu s}\partial^2_2h(X_\mu,X_s)\frac{\partial ({\bf K}_0^{-1})_{\ell m}}{\partial Q_{\mu s}} 
\Biggr\} \delta x_\ell \delta x_m ,\nonumber
\end{eqnarray}
and its average is given by
\begin{eqnarray}
   \langle\mathscr{F}_{\alpha\beta}\rangle  &=&-\frac{1}{2}\biggl\{ \text{Tr}\left( \frac{\partial {\bf K}_0^{-1}}{\partial Q_{\alpha \alpha}}{\bf K}_0\right) \partial_1 h(X_\alpha, X_\beta)+ \text{Tr}\left(\frac{\partial {\bf K}_0^{-1}}{\partial Q_{\alpha \beta}}{\bf K}_0\right)\partial_2h(X_\alpha, X_\beta)\biggr\} \nonumber\\
    & &+\frac{h(X_\alpha, X_\beta)}{2}\sum_s Q_{s\alpha}^{-1}  \Biggl\{
 \left[f_s''+\sum_n W_{sn}\partial_1^2h(X_s,X_n)\right] \text{Tr}\left(  \frac{\partial {\bf K}_0^{-1}}{\partial Q_{ss}}  {\bf K}_0\right) \label{avescrFab}\\
& &+ \sum_\mu W_{\mu s}\partial_1\partial_2h(X_\mu,X_s)\text{Tr}\left(  \frac{\partial {\bf K}_0^{-1}}{\partial Q_{\mu\mu}}  {\bf K}_0\right)+ \nonumber \\
 & &+ \sum_\nu  \text{Tr}\left(  \frac{\partial {\bf K}_0^{-1}}{\partial Q_{s\nu}}  {\bf K}_0\right) W_{s\nu}\partial_1\partial_2h(X_s,X_\nu) + \sum_\mu  \text{Tr}\left(  \frac{\partial {\bf K}_0^{-1}}{\partial Q_{\mu s}}  {\bf K}_0\right) W_{\mu s}\partial^2_2h(X_\mu,X_s)  
   \Biggr\}\nonumber .
\end{eqnarray}
Other relevant averages can be calculated as well:
\begin{equation}
\langle \delta x_i {\mathscr F}_{\alpha\beta}\rangle= \frac{\partial X_i}{\partial W_{\alpha\beta}} \label{dxscrFab},
\end{equation} 
\begin{eqnarray}
    \langle \delta x_i  \delta x_j{\mathscr F}_{\alpha\beta}\rangle &=&\frac{1}{2}\Bigg\{ \partial_1h(X_\alpha,X_\beta)\Bigg[ ({\bf K}_0)_{ij}\text{Tr}\Bigg({\bf K}_0^{-1} \frac{\partial {\bf K}_0}{\partial Q_{\alpha\alpha}} \Bigg) + 2\Bigg( \frac{\partial {\bf K}_0}{\partial Q_{\alpha\alpha}} \Bigg)_{ij} \Bigg]\\
    &+&\partial_2h(X_\alpha,X_\beta)\Bigg[ ({\bf K}_0)_{ij}\text{Tr}\Bigg({\bf K}_0^{-1} \frac{\partial {\bf K}_0}{\partial Q_{\alpha\beta}} \Bigg) + 2\Bigg( \frac{\partial {\bf K}_0}{\partial Q_{\alpha\beta}} \Bigg)_{ij} \Bigg]\Bigg\}\nonumber\\
    &-&\frac{1}{2}h(X_\alpha,X_\beta)\sum_s Q^{-1}_{s\alpha}\Bigg\{ f_s''\Bigg[ ({\bf K}_0)_{ij}\text{Tr}\Bigg({\bf K}_0^{-1}\frac{\partial {\bf K}_0}{\partial Q_{ss}}\Bigg)+2\Bigg( \frac{\partial {\bf K}_0}{\partial Q_{ss}} \Bigg)_{ij} \Bigg]\nonumber\\
    &+&\sum_\mu W_{\mu s}\partial_1\partial_2 h(X_\mu, X_s)\Bigg[ ({\bf K}_0)_{ij}\text{Tr}\Bigg({\bf K}_0^{-1}\frac{\partial {\bf K}_0}{\partial Q_{\mu\mu}}\Bigg)+2\Bigg( \frac{\partial {\bf K}_0}{\partial Q_{\mu\mu}} \Bigg)_{ij} \Bigg]\nonumber\\
    &+&\sum_n W_{sn}\partial_1^2 h(X_s, X_n)\Bigg[ ({\bf K}_0)_{ij}\text{Tr}\Bigg({\bf K}_0^{-1}\frac{\partial {\bf K}_0}{\partial Q_{ss}}\Bigg)+2\Bigg( \frac{\partial {\bf K}_0}{\partial Q_{ss}} \Bigg)_{ij} \Bigg]\nonumber\\
    &+&\sum_\nu W_{s\nu}\partial_1\partial_2 h(X_s, X_\nu)\Bigg[ ({\bf K}_0)_{ij}\text{Tr}\Bigg({\bf K}_0^{-1}\frac{\partial {\bf K}_0}{\partial Q_{s\nu}}\Bigg)+2\Bigg( \frac{\partial {\bf K}_0}{\partial Q_{s\nu}} \Bigg)_{ij} \Bigg]\nonumber\\
    &+&\sum_\mu W_{\mu s}\partial_2^2 h(X_\mu, X_s)\Bigg[ ({\bf K}_0)_{ij}\text{Tr}\Bigg({\bf K}_0^{-1}\frac{\partial {\bf K}_0}{\partial Q_{\mu s}}\Bigg)+2\Bigg( \frac{\partial {\bf K}_0}{\partial Q_{\mu s}} \Bigg)_{ij} \Bigg]\Bigg\}.\nonumber
\end{eqnarray}
For undirected networks, $W_{\alpha\beta}=W_{\beta\alpha}$, so changing the connection from node $\beta$ to node $\alpha$ also changes the connection from node $\alpha$ to node $\beta$ by the same amount. Hence the response function for an undirected network under a change of $W_{\alpha\beta}$ is simply the sum of the two directed-network response functions associated with the changes of $W_{\alpha\beta}$ and $W_{\beta\alpha}$.

\subsection{Conjugate force with respect to the change in the noise strength $\sigma$}
\quad Unlike perturbations of deterministic parameters, variations in the noise intensity do not shift the fixed point, but instead directly modify the steady-state fluctuations. Because the deterministic force does not depend on $\sigma$, we have
\begin{equation}
    ({\bf Q}^{-1}\nabla_\sigma \vec{F}\big|_{\vec{X}})_{s\alpha} = 0.
\end{equation}
The conjugate force associated with a change in the noise strength is then
\begin{eqnarray}
    \mathscr{F}_{\alpha\alpha} &=& -\frac{\partial \Phi}{\partial \sigma_{\alpha\alpha}} = -\frac{1}{2} \sum_{\ell, m} \frac{\partial ({\bf K}_0^{-1})_{\ell m}}{\partial \sigma_{\alpha\alpha}} \delta x_\ell \delta x_m,\\
    \langle \mathscr{F}_{\alpha\alpha} \rangle &=& \frac{1}{2} \text{Tr} \left( {\bf K}_0^{-1} \frac{\partial {\bf K}_0}{\partial \sigma_{\alpha\alpha}} \right).
\end{eqnarray}
The fixed point does not change when the noise strengths are perturbed:
\begin{equation}
    \langle \delta x_i \mathscr{F}_{\alpha\alpha}\rangle = \frac{\partial X_i}{\partial \sigma_{\alpha\alpha}}=0.
\end{equation}
The two-point correlation is
\begin{equation}
    \langle \delta x_i \delta x_j \mathscr{F}_{\alpha\alpha} \rangle = \frac{1}{2} \left[ (\mathbf{K}_0)_{ij} \text{Tr}\left( \mathbf{K}_0^{-1}\frac{\partial \mathbf{K}_0}{\partial \sigma_{\alpha\alpha}} \right) + 2 \left( \frac{\partial \mathbf{K}_0}{\partial \sigma_{\alpha\alpha}} \right)_{ij} \right]
\end{equation}

\subsection{General formula of the time-dependent conjugate force}
\quad For time-dependent perturbations, one needs the time-lagged correlation function between the observable and the conjugate force. To calculate the time-dependent change of the fixed point, in analogy with (\ref{scriptF}), we consider $\langle \delta \vec{x}(\tau) \vec{\mathscr{F}}^\intercal(0) \rangle$. Here $\tau$ denotes the time lag appearing in the correlation functions, while $t$ is reserved for the time variable in the response functions. For the linearized dynamics,
\begin{equation}
    \delta \vec{x}(\tau)=e^{\tau \mathbf{Q}} \delta \vec{x}(0)+ \int_0^\tau ds\, e^{(\tau-s)\mathbf{Q}}\vec{\eta}(s).
\end{equation}
Therefore
\begin{eqnarray}
    \langle \delta x_i(\tau) {\mathscr F}_\alpha(0)\rangle &=& \sum_{\ell} \left( e^{\tau{\bf Q}} \right)_{i\ell}\langle \delta x_\ell(0) {\mathscr F}_\alpha(0)\rangle \nonumber\\
    &+&\int_0^\tau ds\, \sum_k \left( e^{(\tau-s){\bf Q}} \right)_{ik}\langle \eta_k(s) {\mathscr F}_\alpha(0)\rangle.
\end{eqnarray}
Since $s>0$ in the noise integral and ${\mathscr F}_\alpha(0)$ depends only on the state at time $0$, the future noise is uncorrelated with ${\mathscr F}_\alpha(0)$, so the second term vanishes. Using the equal-time identity $\langle \delta x_\ell {\mathscr F}_\alpha \rangle=\partial X_\ell/\partial p_\alpha$, one obtains
\begin{equation}
    \langle \delta x_i(\tau) {\mathscr F}_\alpha(0)\rangle= \sum_{\ell} \frac{\partial X_\ell}{\partial p_\alpha}
\left( e^{\tau{\bf Q}} \right)_{i\ell}, \label{dxtauscrF}
\end{equation}
\begin{eqnarray}
\langle \delta x_i(\tau) \delta x_j(\tau){\mathscr F}_\alpha(0)\rangle&=&  \Biggl\{- \frac{1}{2} \sum_{m,\ell} \sum_{\mu,\nu} \frac{\partial ({\bf K}_0^{-1})_{m\ell}} {\partial Q_{\mu\nu}}\frac {\partial Q_{\mu\nu}}{\partial p_\alpha}\\
& &+\frac{1}{2}  \sum_s \left({\bf Q}^{-1} \nabla_p{\vec { F}} |_{\vec X}\right)_{s\alpha} \sum_{m,\ell} \Bigg[ [f_s''+\sum_n W_{sn}  \partial^2_1 h(X_s,X_n)] \frac{\partial ({\bf K}_0^{-1})_{m\ell}} {\partial Q_{ss}}\nonumber\\
 & & + \sum_\mu W_{\mu s} \partial_1\partial_2h(X_\mu,X_s)\frac{\partial ({\bf K}_0^{-1})_{m\ell}} {\partial Q_{\mu\mu}}+ \sum_\nu W_{s\nu}\partial_1\partial_2h(X_s,X_\nu)\frac{\partial ({\bf K}_0^{-1})_{m\ell}}{\partial Q_{s\nu}}  \nonumber \\
 & &+ \sum_\mu W_{\mu s}\partial^2_2h(X_\mu,X_s) \frac{\partial ({\bf K}_0^{-1})_{m\ell}}{\partial Q_{\mu s}} \Bigg]
\Biggr\} \langle \delta x_i(\tau) \delta x_j(\tau)    \delta x_\ell(0) \delta x_m(0) \rangle \nonumber\\
  &=&- \frac{1}{2}  \sum_{\mu,\nu}\frac {\partial Q_{\mu\nu}}{\partial p_\alpha}\biggl[ ({\bf K}_0)_{ij}\text{Tr}\left(\frac{\partial {\bf K}_0^{-1}}{\partial Q_{\mu \nu}} {\bf K}_0\right)+2 \left({\bf K}_\tau\frac{\partial {\bf K}_0^{-1}}{\partial Q_{\mu \nu}} {\bf K}_\tau^\intercal\right)_{ij}\bigg]\nonumber\\
  &&+\frac{1}{2}  \sum_s \left({\bf Q}^{-1} \nabla_p{\vec { F}} |_{\vec X}\right)_{s\alpha}  \Biggl\{ [ f_s''+  \sum_n W_{sn}\partial_1^2h(X_s,X_n)] \times\nonumber\\
    & &\biggl[ ({\bf K}_0)_{ij}\text{Tr}\left(\frac{\partial {\bf K}_0^{-1}}{\partial Q_{ss}} {\bf K}_0\right)+2 \left({\bf K}_\tau\frac{\partial {\bf K}_0^{-1}}{\partial Q_{ss}} {\bf K}_\tau^\intercal\right)_{ij}\bigg] \label{dxdxtauscrF}\\
   & & +\sum_\mu W_{\mu s}\partial_1\partial_2h(X_\mu,X_s) \biggl[ ({\bf K}_0)_{ij}\text{Tr}\left(\frac{\partial {\bf K}_0^{-1}}{\partial Q_{\mu\mu}} {\bf K}_0\right)+2 \left({\bf K}_\tau\frac{\partial {\bf K}_0^{-1}}{\partial Q_{\mu\mu}} {\bf K}_\tau^\intercal\right)_{ij}\bigg] \nonumber\\
   & & +  \sum_\nu W_{s\nu} \partial_1\partial_2h(X_s,X_\nu)  \biggl[ ({\bf K}_0)_{ij}\text{Tr}\left(\frac{\partial {\bf K}_0^{-1}}{\partial Q_{s\nu}} {\bf K}_0\right)+2 \left({\bf K}_\tau\frac{\partial {\bf K}_0^{-1}}{\partial Q_{s\nu}} {\bf K}_\tau^\intercal\right)_{ij}\bigg]    \nonumber\\
   & & +\sum_\mu W_{\mu s}  \partial_2^2h(X_\mu,X_s)\biggl[ ({\bf K}_0)_{ij}\text{Tr}\left(\frac{\partial {\bf K}_0^{-1}}{\partial Q_{\mu s}} {\bf K}_0\right)+2 \left({\bf K}_\tau\frac{\partial {\bf K}_0^{-1}}{\partial Q_{\mu s}} {\bf K}_\tau^\intercal\right)_{ij}\bigg]
   \Biggr\} \nonumber
\end{eqnarray}
where 
\begin{eqnarray}
\langle \delta x_i(\tau) \delta x_j (\tau)  \delta x_\ell(0) \delta x_m (0)\rangle
&=&({\bf K}_0)_{ij}({\bf K}_0)_{\ell m} + ({\bf K}_\tau)_{im}({\bf K}_\tau)_{j\ell }+ ({\bf K}_\tau)_{i\ell}({\bf K}_\tau)_{jm} \label{dx4tau}
\end{eqnarray}
which is used in the last equality. Note that ${\bf K}_\tau= e^{\tau{\bf Q}}{\bf K}_0$, and the right-hand side of (\ref{dxdxtauscrF}) can therefore be calculated explicitly in terms of the network model parameters.
The time-dependent response function of the state is
\begin{eqnarray}
    \chi_{i\alpha}(t) &=& \chi_{i\alpha}^{(ss)}-\langle \delta x_i(t) \mathscr{F}_\alpha(0) \rangle\\
    &=& \frac{\partial X_i}{\partial p_\alpha}
-\sum_{\ell} \frac{\partial X_\ell}{\partial p_\alpha}
\left( e^{t{\bf Q}} \right)_{i\ell}\nonumber
\end{eqnarray}
The non-steady response function of the force is
\begin{equation}
    \chi_{i\alpha}(t)= \chi_{i\alpha}^{(ss)}-\langle \delta F_i(t) {\mathscr F}_\alpha(0)\rangle
\end{equation}
\begin{eqnarray}
\langle \delta {F}_i(t) \mathscr F_\alpha(0)\rangle&=& \langle \delta f_i(t)  {\mathscr F}_\alpha(0)\rangle +\sum_{j\neq i} W_{ij} \Biggl\{ \partial_1 h(X_i,X_j) \langle \delta x_i(t) {\mathscr F}_\alpha(0)\rangle \nonumber\\ &+&\partial_2h(X_i,X_j)\langle \delta x_j(t){\mathscr F}_\alpha(0)\rangle 
+\frac{1}{2} \Big[\partial_1^2h(X_i,X_j) \langle \delta x_i^2(t) {\mathscr F}_\alpha(0) \rangle\nonumber\\
&+& \partial_2^2h(X_i,X_j) \langle \delta x_j^2(t) {\mathscr F}_\alpha(0)\rangle \Big] +\partial_1 \partial_2h(X_i,X_j)\langle \delta x_i(t) \delta x_j(t) {\mathscr F}_\alpha(0)\rangle \Biggr\}.
\end{eqnarray}
If we only consider the linear term, then it becomes
\begin{equation}
    \langle \delta F_i(t) \mathscr{F}_\alpha(0) \rangle = \sum_{\ell} ({\bf Q}e^{t{\bf Q}})_{i\ell} \frac{\partial X_\ell}{\partial p_\alpha}.
\end{equation}
The non-steady, time-dependent response function of the fluctuation correlation can then be written as
\begin{equation}
    \chi_{ij\alpha}(t)= \chi_{ij\alpha}^{(ss)}-[\langle \delta x_i(t) \delta x_j(t){\mathscr F}_\alpha(0)\rangle -({\bf K}_0)_{ij}\langle{\mathscr F}_\alpha\rangle ],\hbox{  where }
\end{equation}
\begin{eqnarray}
\langle \delta x_i(t) \delta x_j(t){\mathscr F}_\alpha(0)\rangle -({\bf K}_0)_{ij}\langle{\mathscr F}_\alpha\rangle&=&  \sum_{\mu,\nu}\frac {\partial Q_{\mu\nu}}{\partial p_\alpha}\left(e^{t{\bf Q}}\frac{\partial {\bf K}_0}{\partial Q_{\mu \nu}}e^{t{\bf Q}^\intercal}\right)_{ij}-  \sum_s \left({\bf Q}^{-1} \nabla_p{\vec { F}}|_{\vec X} \right)_{s\alpha} \times\nonumber\\
& &\Biggl\{ [ f_s''+  \sum_n W_{sn}\partial_1^2h(X_s,X_n)] \left(e^{t{\bf Q}}\frac{\partial {\bf K}_0}{\partial Q_{ss}}e^{t{\bf Q}^\intercal}\right)_{ij} \\
   & & +\sum_\mu W_{\mu s}\partial_1\partial_2h(X_\mu,X_s) \left(e^{t{\bf Q}}\frac{\partial {\bf K}_0}{\partial Q_{\mu\mu}}e^{t{\bf Q}^\intercal}\right)_{ij}\nonumber\\
   & & +  \sum_\nu W_{s\nu} \partial_1\partial_2h(X_s,X_\nu)    \left(e^{t{\bf Q}}\frac{\partial {\bf K}_0}{\partial Q_{s\nu}}e^{t{\bf Q}^\intercal}\right)_{ij}    \nonumber\\
   & & +\sum_\mu W_{\mu s}  \partial_2^2h(X_\mu,X_s)\left(e^{t{\bf Q}}\frac{\partial {\bf K}_0}{\partial Q_{\mu s}}e^{t{\bf Q}^\intercal}\right)_{ij}
   \Biggr\}. \nonumber
\end{eqnarray}
